% Source: Weinberg GR, physical PDF pp. 149--151
%         (printed pp. 124--126).
% Coverage: Section 5.2, Electrodynamics; equations (5.2.1)--(5.2.14).
% Figures/tables/footnotes: none.
% Status: source-reviewed and compile-clean.
% Uncertainties: none.

\subsection{Electrodynamics}\label{sec:5.2}

We recall that in the absence of gravitational fields the Maxwell equations
of electrodynamics can be written as
\begin{equation}
  \partial_a F^{ab}=-J^b,
  \tag{5.2.1}\label{eq:5.2.1}
\end{equation}
\begin{equation}
  \partial_{[a}F_{bc]}=0,
  \tag{5.2.2}\label{eq:5.2.2}
\end{equation}
where \(J^a=(\rho,\mathbf J)\) is the current four-vector and \(F^{ab}\) is
the field-strength tensor, with
\(F^{0i}=E_i,\ F^{ij}=\varepsilon^{ijk}B_k\), and so on. (See
Section~\ref{sec:2.7}.) Suppose that we define \(F^{\mu\nu}\) and \(J^\mu\)
in general coordinates by the requirements that they reduce to \(F^{ab}\)
and \(J^a\) in locally inertial Minkowskian coordinates, and that they behave
as tensors under general coordinate transformations. That is, if \(F^{ab}\)
and \(J^a\) are the values measured in a locally inertial frame, then
\[
  F^{\mu\nu}
  \equiv
  \frac{\partial x^\mu}{\partial\xi^a}
  \frac{\partial x^\nu}{\partial\xi^b}F^{ab},
  \qquad
  J^\mu
  \equiv
  \frac{\partial x^\mu}{\partial\xi^a}J^a.
\]
We can then make Eqs.~\eqref{eq:5.2.1} and \eqref{eq:5.2.2} generally
covariant by replacing all derivatives by covariant derivatives:
\begin{equation}
  \nabla_\mu F^{\mu\nu}=-J^\nu,
  \tag{5.2.3}\label{eq:5.2.3}
\end{equation}
\begin{equation}
  \nabla_{[\mu}F_{\nu\rho]}=0,
  \tag{5.2.4}\label{eq:5.2.4}
\end{equation}
it being now understood that indices are to be raised and lowered with
\(g_{\mu\nu}\) instead of \(\eta_{\mu\nu}\); that is,
\begin{equation}
  F_{\mu\nu}
  \equiv
  g_{\mu\rho}g_{\nu\sigma}F^{\rho\sigma}.
  \tag{5.2.5}\label{eq:5.2.5}
\end{equation}

Since \(F^{\mu\nu}\) and \(F_{\mu\nu}\) are antisymmetric, we may use
Eqs.~\eqref{eq:4.7.10} and \eqref{eq:4.7.11} to rewrite the Maxwell
equations as
\begin{equation}
  \partial_\mu\!\left(\sqrt{-g}\,F^{\mu\nu}\right)
  =
  -\sqrt{-g}\,J^\nu,
  \tag{5.2.6}\label{eq:5.2.6}
\end{equation}
\begin{equation}
  \partial_{[\mu}F_{\nu\rho]}=0.
  \tag{5.2.7}\label{eq:5.2.7}
\end{equation}
Equations~\eqref{eq:5.2.3}--\eqref{eq:5.2.7} are true in the absence of
gravitation and generally covariant and hence, according to the Principle of
General Covariance, also true in arbitrary gravitational fields.

The electromagnetic force on a particle of charge \(e\) is given in the
absence of gravitation by Eq.~\eqref{eq:2.7.9}:
\begin{equation}
  f^a=eF^a{}_b\frac{\dd\xi^b}{\dd\tau}.
  \tag{5.2.8}\label{eq:5.2.8}
\end{equation}
We immediately conclude that in general coordinates the electromagnetic
force in an arbitrary gravitational field is
\begin{equation}
  f^\mu=eF^\mu{}_\nu\frac{\dd x^\nu}{\dd\tau},
  \tag{5.2.9}\label{eq:5.2.9}
\end{equation}
where of course
\begin{equation}
  F^\mu{}_\nu
  \equiv
  g_{\nu\lambda}F^{\mu\lambda}.
  \tag{5.2.10}\label{eq:5.2.10}
\end{equation}
Once again, we are using the Principle of General Covariance;
Eq.~\eqref{eq:5.2.9} obviously reduces to Eq.~\eqref{eq:5.2.8} in locally
inertial Minkowskian coordinates, and it is generally covariant, because
\(f^\mu\) is a vector (Section~\ref{sec:5.1}),
\(\dd x^\nu/\dd\tau\) is a vector, and \(F^\mu{}_\nu\) is defined as a
tensor; therefore Eq.~\eqref{eq:5.2.9} is true.

It is instructive to evaluate the current vector \(J^\mu\). In special
relativity it is
\begin{equation}
  J^a
  =
  \sum_n e_n
  \int
  \delta^4\!\left(\xi-\xi_n\right)\dd\xi_n^a,
  \tag{5.2.11}\label{eq:5.2.11}
\end{equation}
the integral being taken along the trajectory of the \(n\)th particle. [See
Eq.~\eqref{eq:2.6.5}.] The four-dimensional delta function in a general
coordinate system is defined by
\begin{equation}
  \int
  \dd^4x\,\Phi(x)\delta^4(x-y)
  =
  \Phi(y).
  \tag{5.2.12}\label{eq:5.2.12}
\end{equation}
Since \(\sqrt{-g}\,\dd^4x\) is a scalar,
\((-g)^{-1/2}\delta^4(x-y)\) must be a scalar, which of course reduces to
the ordinary delta function in special relativity, where \(-g=1\). (In some
works it is this scalar that is defined as the delta function.) Thus the
contravariant vector that reduces to \(J^a\) in the absence of gravitation
is
\begin{equation}
  J^\mu(x)
  =
  \frac{1}{\sqrt{-g(x)}}
  \sum_n e_n
  \int
  \delta^4\!\left(x-x_n\right)\dd x_n^\mu.
  \tag{5.2.13}\label{eq:5.2.13}
\end{equation}

Note that the conservation law \(\partial_aJ^a=0\) of special relativity
becomes in general relativity \(\nabla_\mu J^\mu=0\), or, using
Eq.~\eqref{eq:4.7.7},
\begin{equation}
  \partial_\mu\!\left(\sqrt{-g}\,J^\mu\right)=0.
  \tag{5.2.14}\label{eq:5.2.14}
\end{equation}
The factor \((-g)^{-1/2}\) in Eq.~\eqref{eq:5.2.13} is just what is needed
to cancel the \(\sqrt{-g}\) in Eq.~\eqref{eq:5.2.14}, so that
Eq.~\eqref{eq:5.2.14} still just expresses the constancy of the \(e_n\).
